% !TEX program = pdflatex
% !TEX encoding = UTF-8 Unicode
% ==============================================================================
% Segundo Relatório / Artigo Científico — Estritamente conforme template_de_prompt_para_artigo.md
% Projeto: OntoML-Mod (Ontologia de Machine Learning) — CIn / UFPE (2026.2)
% Base Ontológica Consolidada: Projeto/Ontology_02.ttl (Ciclos 1, 2 e 3)
% ==============================================================================
\documentclass[11pt,a4paper]{article}

\usepackage[utf8]{inputenc}
\usepackage[T1]{fontenc}
\usepackage[english,brazil]{babel}
\usepackage{lmodern}
\usepackage{microtype}
\usepackage[a4paper,top=2.5cm,bottom=2.5cm,left=2.5cm,right=2.5cm]{geometry}
\usepackage{setspace}
\onehalfspacing

\usepackage{amsmath,amssymb}
\usepackage{booktabs}
\usepackage{tabularx}
\usepackage{array}
\usepackage{enumitem}
\usepackage{url}
\usepackage{xcolor}
\usepackage{titlesec}
\usepackage[hidelinks]{hyperref}

% Previne estouro de margem em termos tecnicos longos e identificadores CamelCase
\emergencystretch=1.8em

% Cores sóbrias institucionais
\definecolor{cinblue}{HTML}{1B365D}
\definecolor{codebg}{HTML}{F8FAFC}

% Macro para termos ontológicos em fonte monoespaçada
\newcommand{\onto}[1]{\texttt{\small #1}}

% Formatação enxuta de títulos para respeitar estritamente os limites de página do template
\titleformat{\section}{\large\bfseries\color{cinblue}}{}{0em}{}
\titlespacing*{\section}{0pt}{2.0ex plus 0.5ex minus .2ex}{1.0ex plus .2ex}
\titleformat{\subsection}{\normalsize\bfseries}{}{0em}{}
\titlespacing*{\subsection}{0pt}{1.4ex plus 0.4ex minus .2ex}{0.6ex plus .2ex}

\setlength{\parindent}{1.25cm}
\setlength{\parskip}{3pt}

% ==============================================================================
\title{\vspace{-1.0cm}\textbf{\Large OntoML-Mod: Uma Ontologia em OWL~2~DL e Lógicas Descritivas para Representação Semântica de Modelos, Fundamentos Matemáticos e Aplicações de Machine Learning}}

\author{
  \textbf{Gabriel M.\ de Albuquerque}$^{1}$, \textbf{Gabriel S.\ Arnaud}$^{1}$\\[4pt]
  \small $^{1}$Centro de Informática --- Universidade Federal de Pernambuco (CIn-UFPE)\\
  \small Caixa Postal 7809 -- 50.732-970 -- Recife -- PE -- Brasil\\
  \small \texttt{\{gma5, gsa6\}@cin.ufpe.br}
}
\date{\small Novembro de 2026}

\begin{document}

\maketitle
\vspace{-0.4cm}

\begin{abstract}
\noindent\textbf{Resumo.} O avanço acelerado do Aprendizado de Máquina (\emph{Machine Learning} -- ML) produziu uma ampla diversidade de famílias algorítmicas, formulações matemáticas e métricas de validação frequentemente documentadas de modo fragmentado na literatura. Este artigo apresenta a \textbf{OntoML-Mod}, uma ontologia de domínio em OWL~2~DL e Lógicas Descritivas desenvolvida no ambiente Protégé segundo a metodologia \emph{Methontology}. Diferentemente de esquemas voltados apenas a registros de execução experimental ou catálogos de bibliotecas, a OntoML-Mod articula cinco eixos conceituais --- \onto{Algorithms}, \onto{MLTypes}, \onto{MathematicalFoundations}, \onto{EvaluationMetrics} e \onto{Applications} --- por meio de propriedades de objeto inversas, propriedades de dados funcionais e axiomas em Lógica Descritiva validados na literatura científica para algoritmos de agrupamento, modelos baseados em árvores e redes neurais artificiais.

\smallskip
\noindent\textbf{Palavras-chave:} Ontologia de Domínio; Aprendizado de Máquina; OWL~2~DL; Lógicas Descritivas; Methontology; Protégé.
\end{abstract}

\begin{otherlanguage}{english}
\begin{abstract}
\noindent\textbf{Abstract.} The rapid growth of Machine Learning (ML) has generated a wide diversity of algorithmic families, mathematical formulations, and validation metrics that often remain scattered across the literature. This paper presents \textbf{OntoML-Mod}, a domain ontology in OWL~2~DL and Description Logics developed in Protégé following the \emph{Methontology} methodology. Unlike schemas restricted to experimental run logs or software library catalogs, OntoML-Mod connects five core axes---\onto{Algorithms}, \onto{MLTypes}, \onto{MathematicalFoundations}, \onto{EvaluationMetrics}, and \onto{Applications}---through inverse object properties, functional datatype properties, and Description Logic axioms grounded in peer-reviewed literature across clustering, tree-based models, and artificial neural networks.

\smallskip
\noindent\textbf{Keywords:} Domain Ontology; Machine Learning; OWL~2~DL; Description Logics; Methontology; Protégé.
\end{abstract}
\end{otherlanguage}

% ==============================================================================
\section{1. INTRODUÇÃO}
% ==============================================================================

Historicamente, o Aprendizado de Máquina (\emph{Machine Learning} -- ML) consolidou-se no âmbito da Inteligência Artificial e da Ciência da Computação como o paradigma voltado à construção de modelos computacionais capazes de extrair padrões, generalizar a partir de dados e aprimorar seu desempenho preditivo ou decisório sem programação explícita de cada regra, tornando-se infraestrutura estratégica em áreas como saúde, finanças, indústria, segurança digital e sistemas ciberfísicos~\cite{sarker2021ml}. Em decorrência dessa expansão acelerada e contínua, entretanto, o conhecimento acumulado da área --- compreendendo distintos regimes de aprendizado, famílias de algoritmos, formulações teóricas e critérios de avaliação --- encontra-se amplamente fragmentado e disperso em publicações heterogêneas~\cite{sarker2021ml}. De modo paralelo e independente, o campo da Engenharia de Ontologias, enraizado na Representação do Conhecimento e na Web Semântica, consolidou métodos rigorosos para estruturar conceitualizações formais, explícitas e compartilhadas de domínios complexos por meio de vocabulários padronizados, hierarquias taxonômicas, relações semânticas e axiomas lógicos~\cite{fernandez1997methontology}. Na atualidade, essas especificações formais desempenham papel central ao conferir organização conceitual, interoperabilidade semântica entre humanos e sistemas computacionais, reutilização de conhecimento e suporte à inferência lógica automatizada~\cite{fernandez1997methontology}.

A interseção entre essas duas disciplinas tem motivado a construção de artefatos semânticos voltados à organização formal do ecossistema de Aprendizado de Máquina, a começar pela \textbf{MLOnto}, proposta por Braga, Dias e Regateiro~\cite{braga2020mlonto}, cujo objetivo é estruturar uma base de conhecimento explícita e inferível sobre componentes do domínio para apoiar tanto múltiplos interessados humanos quanto agentes autônomos de software. De forma complementar, a \textbf{ML Ontology} desenvolvida por Humm~\cite{humm2026industry} tem como propósito atuar como grafo de conhecimento leve voltado à prontidão industrial, alinhando terminologias de bibliotecas e ferramentas de automação para sustentar assistentes de configuração e ajuda interativa em plataformas de ciência de dados. Em outra vertente, El~Massari et al.~\cite{elmassari2022diabetes} buscam integrar regras extraídas de modelos preditivos em ontologias clínicas com raciocínio lógico para ampliar a interpretabilidade e o desempenho no apoio ao diagnóstico médico. Por fim, o \textbf{ML-Schema}, concebido por Publio et al.~\cite{publio2018mlschema}, objetiva fornecer um esquema canônico para padronizar o intercâmbio e a reprodutibilidade de metadados experimentais entre plataformas heterogêneas, harmonizando iniciativas correlatas voltadas à anotação de entidades de mineração de dados (\emph{OntoDM-core}), à otimização e meta-mineração de fluxos (\emph{DMOP}) e ao compartilhamento de experimentos (\emph{Exposé} e vocabulário \emph{MEX}). Contudo, embora atendam com êxito à atuação de agentes autônomos, à operação industrial, ao diagnóstico clínico ou à rastreabilidade experimental, essas iniciativas não têm como foco central atuar como um guia estruturado de aprendizagem e estudo para a área, carecendo de uma visão conceitual integrada que relacione didaticamente paradigmas de aprendizado, modelos, fundamentos matemáticos, métricas de avaliação e aplicações.

Essa lacuna torna-se ainda mais crítica diante da velocidade com que novas metodologias emergem na literatura, fazendo com que estudantes, pesquisadores e profissionais em formação frequentemente empreguem modelos computacionais como ``caixas-pretas'' parametrizadas, sem compreensão clara dos pressupostos teóricos que fundamentam seu funcionamento e condicionam sua escolha~\cite{sarker2021ml}. Sob uma perspectiva formativa e científica, o domínio sólido do Aprendizado de Máquina requer articular em um mesmo arcabouço conceitual como diferentes paradigmas se traduzem em famílias algorítmicas, quais teorias matemáticas os sustentam, por quais critérios seus resultados devem ser avaliados e em que contextos reais são aplicáveis. Justifica-se, portanto, a concepção de uma ontologia construída sob princípios sistemáticos da Engenharia Ontológica~\cite{fernandez1997methontology}, capaz de operar simultaneamente como representação formal decidível em lógica descritiva e como compêndio conceitual navegável de referência para estudo e consolidação do conhecimento na disciplina.

Diante desse contexto, o objetivo deste trabalho é desenvolver, formalizar e validar a \textbf{OntoML-Mod}, uma ontologia de domínio em OWL~2~DL fundamentada em literatura científica revisada por pares e orientada por Questões de Competência~\cite{bezerra2013evaluating,potoniec2020dataset}, concebida como um guia estruturado de aprendizagem e referência conceitual para o Aprendizado de Máquina. Ao organizar semanticamente e interconectar paradigmas de aprendizado, categorias de modelos, alicerces matemáticos, critérios de avaliação, características operacionais e domínios de aplicação em uma estrutura formal consistente e consultável, a ontologia proposta busca apoiar tanto o estudo sistemático e a apropriação teórica da área por iniciantes e especialistas quanto a descoberta automatizada de conhecimento por raciocinadores e consultas semânticas.

% ==============================================================================
\section{2. TRABALHOS RELACIONADOS}
% ==============================================================================

\subsection{2.1. Machine Learning Ontology (MLOnto)}

Proposta por Braga, Dias e Regateiro~\cite{braga2020mlonto}, a \textbf{MLOnto} constitui a principal referência de base para o escopo conceitual deste projeto. Desenvolvida no Protégé~5 em OWL com adoção da sintaxe Manchester, a MLOnto foi concebida para fornecer um repositório de conhecimento explícito e inferível voltado tanto a múltiplos \emph{stakeholders} humanos quanto a agentes autônomos de software integrantes do modelo \emph{Autonomous Architecture Over Restricted Domains} (A2RD) no ambiente SKAU~\cite{braga2020mlonto}. Sua arquitetura estrutura o domínio de Aprendizado de Máquina em sete superclasses de topo: \onto{Algorithms}, \onto{Applications}, \onto{Dependencies}, \onto{Dictionary}, \onto{Frameworks}, \onto{Involved} e \onto{MLTypes} (esta última subdividida em \onto{AutoML}, \onto{Reinforcement\_Learning}, \onto{Semi-supervised\_Learning}, \onto{Supervised\_Learning} e \onto{Unsupervised\_Learning}). Ao todo, a MLOnto contabiliza 433 classes, 56 propriedades de objeto, 1 propriedade de dados, 1 indivíduo, 11 propriedades de anotação e 1.877 axiomas~\cite{braga2020mlonto}.

\subsection{2.2. ML Ontology (An Industry-Ready Machine Learning Ontology)}

Desenvolvida por Humm~\cite{humm2026industry} ao longo de mais de cinco anos com a colaboração de mais de 15 engenheiros do conhecimento, a \textbf{ML Ontology} tem como foco central a prontidão industrial (\emph{industry-readiness}), atuando como espinha dorsal informacional da plataforma de ciência de dados e \emph{Meta-AutoML} denominada OMA-ML. O artefato sustenta dois casos de uso principais: um assistente de configuração de treinamento capaz de harmonizar parâmetros entre bibliotecas e soluções heterogêneas (como \emph{Scikit-learn}, \emph{Auto-Sklearn} e \emph{AutoKeras}) e um sistema de ajuda interativa com cerca de 350 termos alinhados ao \emph{Wikidata}~\cite{humm2026industry}. Para privilegiar simplicidade operacional e desempenho de consulta (inferior a 100\,ms), adota uma modelagem leve em RDF/RDFS, SKOS e SHACL/SPARQL composta por cerca de 10 classes organizadas nos módulos \emph{ML concepts} e \emph{ML implementations}, com aproximadamente 700 indivíduos e 5.000 triplas RDF, obtendo F1-score de 0,83 em avaliação terminológica (precisão de 1,00 e revocação de 0,71, caindo para 0,38 na classe de áreas de ML pela omissão inicial de paradigmas como aprendizado semi-supervisionado, auto-supervisionado e por transferência)~\cite{humm2026industry}.

\subsection{2.3. Ontologia para Predição de Diabetes com Machine Learning (El Massari et al.)}

El~Massari et al.~\cite{elmassari2022diabetes} apresentam uma ontologia aplicada ao diagnóstico clínico de diabetes que combina seis algoritmos clássicos de classificação supervisionada (\emph{SVM}, \emph{KNN}, \emph{ANN}, \emph{Naive Bayes}, Regressão Logística e Árvore de Decisão) com raciocínio ontológico em OWL e regras SWRL sobre o conjunto \emph{Pima Indians Diabetes Database} (768 instâncias). No Protégé, os autores modelam classes de domínio clínico (\onto{Patient}, \onto{Diagnostic}, \onto{PatientEvaluationMetrics} com \onto{TP}, \onto{TN}, \onto{FP} e \onto{FN}), importam os registros via plugin \emph{Cellfie} e convertem as regras de decisão induzidas por Árvore de Decisão no Weka em regras lógicas SWRL (\emph{SWRLTab}) processadas pelo raciocinador \emph{Pellet}. A contribuição central reside em demonstrar que o classificador ontológico interpretável atinge acurácia superior (77,5\% em validação cruzada de 10 pastas) e alta precisão (81,2\%) frente aos classificadores isolados avaliados por \onto{Accuracy}, \onto{Precision}, \onto{Recall}, \onto{F-Measure} e \onto{ROC Area}~\cite{elmassari2022diabetes}.

\subsection{2.4. ML-Schema e Ontologias de Fluxos Experimentais (Publio et al.)}

Proposto por Publio et al.~\cite{publio2018mlschema} no âmbito do \emph{W3C Machine Learning Schema Community Group}, o \textbf{ML-Schema} é um esquema ontológico canônico composto por aproximadamente 25 classes de alto nível e projetado para padronizar a representação e o intercâmbio de metadados sobre experimentos de Aprendizado de Máquina, mitigando problemas de interoperabilidade e reprodutibilidade entre repositórios como o \emph{OpenML}. Suas classes nucleares concentram-se na rastreabilidade do ciclo experimental --- incluindo \onto{Study}, \onto{Experiment}, \onto{Run}, \onto{Algorithm}, \onto{Implementation}, \onto{HyperParameter}, \onto{Data}, \onto{Dataset}, \onto{Task}, \onto{Model}, \onto{ModelEvaluation} e \onto{EvaluationMeasure} --- além de fornecer mapeamentos diretos para ontologias e vocabulários precedentes de mineração de dados: \textbf{OntoDM-core} (alinhada a OBI/BFO em três camadas de especificação, implementação e aplicação), \textbf{DMOP} (focada em meta-mineração e otimização em ``caixa de vidro'' de \emph{workflows}), \textbf{Exposé} e o vocabulário leve \textbf{MEX}~\cite{publio2018mlschema}.

\subsection{2.5. Lacuna Científica}

Embora os trabalhos relacionados ofereçam contribuições consolidadas, persistem lacunas relevantes para a representação conceitual e didático-científica do Aprendizado de Máquina: (i) a \textbf{MLOnto}~\cite{braga2020mlonto} não dispõe de superclasses nem restrições dedicadas a fundamentos matemáticos ou métricas de avaliação, acopla algoritmos diretamente sob paradigmas de aprendizado e omite abordagens contemporâneas como \emph{Self-Supervised Learning}~\cite{liu2023self} e \emph{Transfer Learning}~\cite{zhuang2021transfer}; (ii) a \textbf{ML Ontology} de Humm~\cite{humm2026industry} prioriza a interoperabilidade leve em RDFS/SKOS voltada à parametrização de bibliotecas e \emph{AutoML}, abdicando da axiomatização em OWL~2~DL sobre teorias matemáticas e apresentando baixa revocação em paradigmas de aprendizado; (iii) o modelo de \textbf{El~Massari et al.}~\cite{elmassari2022diabetes} formaliza entidades clínicas de pacientes para inferência diagnóstica específica via SWRL, não constituindo uma ontologia conceitual sobre modelos de ML; e (iv) o \textbf{ML-Schema} e as ontologias mapeadas por Publio et al.~\cite{publio2018mlschema} privilegiam o registro de execuções experimentais e a otimização de fluxos de mineração de dados, sem constituir um mapa conceitual enxuto orientado à aprendizagem. A \textbf{OntoML-Mod} preenche essa lacuna ao formalizar em OWL~2~DL uma ontologia densamente axiomatizada e validada bibliograficamente que atua como guia estruturado de estudo, vinculando famílias algorítmicas contemporâneas aos seus fundamentos matemáticos, paradigmas de supervisão, métricas, propriedades operacionais e aplicações.

% ==============================================================================
\section{3. MÉTODOS}
% ==============================================================================

\subsection{3.1. Planejamento e Escopo}

A construção da \textbf{OntoML-Mod} fundamentou-se na metodologia \textbf{Methontology}~\cite{fernandez1997methontology}, amplamente adotada na Engenharia Ontológica por estruturar o ciclo de vida por protótipos evolutivos através das atividades de especificação, aquisição de conhecimento, conceitualização, formalização, implementação, avaliação e documentação. Dessa forma, a aplicação da Methontology proporcionou maior rigor metodológico e semântico durante a construção da ontologia. Diferentemente de ontologias exaustivas de infraestrutura, o escopo foi planejado para modelar os tipos de algoritmos de ML em nível conceitual e estrutural, conectando-os às suas aplicações práticas, métricas e bases matemáticas. Para assegurar evolução controlada e rastreabilidade bibliográfica estrita ao longo do desenvolvimento no Protégé, a equipe estabeleceu e executou o seguinte \emph{pipeline} geral iterativo em quatro etapas:
\begin{enumerate}[leftmargin=*,nosep]
  \item \textbf{Escolher e focar em uma subclasse de algoritmos:} delimitação modular da subárvore-alvo a cada ciclo evolutivo (Ciclo~1 sobre \onto{MLTypes} e \onto{ClusteringAlgorithms}; Ciclo~2 sobre propriedades operacionais e a subárvore \onto{TreeBased}; e Ciclo~3 sobre a subárvore \onto{ArtificialNeuralNetwork}, abrangendo arquiteturas rasas, profundas e autocodificadores);
  \item \textbf{Pesquisar artigos científicos que embasarão os assuntos matemáticos, aplicações, etc.:} seleção criteriosa de obras seminais e revisões sistemáticas para validar cada classe, restrição matemática, métrica e domínio de aplicação, podando conceitos sem respaldo nas fontes selecionadas;
  \item \textbf{Atualizar o Protégé:} codificação e saneamento de classes, propriedades de objeto e de dados, axiomas de disjunção, restrições semânticas e indivíduos na ontologia;
  \item \textbf{Atualizar o relatório:} documentação formal das decisões de modelagem, métricas estruturais e resultados de validação.
\end{enumerate}

\subsection{3.2. Aquisição do Conhecimento}

A aquisição do conhecimento ocorreu de forma contínua ao longo dos três ciclos do \emph{pipeline} e baseou-se em artigos científicos revisados por pares e obras clássicas extraídos predominantemente do \textbf{Google Scholar}, além da análise estrutural das ontologias correlatas~\cite{publio2018mlschema,braga2020mlonto,elmassari2022diabetes,humm2026industry}. Para fundamentar a divisão de \onto{MLTypes} em paradigmas disjuntos no Ciclo~1, utilizaram-se três artigos-base: Sarker~\cite{sarker2021ml}, Zhuang et al.~\cite{zhuang2021transfer} e Liu et al.~\cite{liu2023self}. Para a subárvore \onto{ClusteringAlgorithms} e suas propriedades de dados intrínsecas, o conhecimento foi extraído de quatro revisões e artigos de referência --- Nagpal, Jatain e Gaur~\cite{nagpal2013review}, Yin et al.~\cite{yin2026rapid}, de~Carvalho et al.~\cite{decarvalho2018gaussian} e Simões e de~Carvalho~\cite{simoes2023gaussian} --- complementados pelas fontes primárias de Gibson, Kleinberg e Raghavan~\cite{gibson1998clustering} para \onto{STIRR} e de Dhillon e Modha~\cite{dhillon2001concept} para \onto{SphericalKMeans}. No Ciclo~2, dedicado aos algoritmos baseados em árvores (\onto{TreeBased}), selecionaram-se seis novas obras de alta densidade conceitual e tradição teórica --- Breiman et al.~\cite{breiman1984cart} (\onto{CART}), Quinlan~\cite{quinlan1993c45} (\onto{ID3} e \onto{C4\_5}), Freund e Schapire~\cite{freund1997decision} (\onto{AdaBoost}), Breiman~\cite{breiman2001random} (\onto{RandomForest}), Geurts, Ernst e Wehenkel~\cite{geurts2006extremely} (\onto{ExtraTrees}) e Bentéjac, Csörgő e Martínez-Muñoz~\cite{bentejac2021comparative} (\onto{GradientBoostingMachine}, \onto{ExtremeGB}, \onto{LightGBMachine} e \onto{CategoricalBoosting}) --- somadas a Sarker~\cite{sarker2021ml} e El~Massari et al.~\cite{elmassari2022diabetes}.

No Ciclo~3, voltado à consolidação completa da subárvore de redes neurais artificiais (\onto{ArtificialNeuralNetwork}) e à expansão correspondente dos fundamentos matemáticos, métricas e aplicações, incorporaram-se quinze referências científicas centrais da literatura de aprendizado profundo e conexionista. A base conceitual geral e as estratégias de regularização apoiaram-se na obra de referência de Goodfellow, Bengio e Courville~\cite{goodfellow2016deep}, articulada aos artigos seminais das arquiteturas clássicas e recorrentes/convolucionais: Rumelhart, Hinton e Williams~\cite{rumelhart1986learning} (retropropagação em redes multicamadas), Broomhead e Lowe~\cite{broomhead1988multivariable} (redes de base radial), LeCun et al.~\cite{lecun1998gradient} (redes neurais convolucionais) e Hochreiter e Schmidhuber~\cite{hochreiter1997lstm} (memória de longo e curto prazo). Para a subfamília especializada de autocodificadores (\onto{AutoEncoders}), utilizaram-se a formulação teórica de Baldi~\cite{baldi2012autoencoders}, os estudos de pré-treinamento em camadas e versões empilhadas de Bengio et al.~\cite{bengio2007greedy} e Vincent et al.~\cite{vincent2010stacked}, a variante de remoção de ruído de Vincent et al.~\cite{vincent2008extracting}, a regularização contrativa de Rifai et al.~\cite{rifai2011contractive}, a esparsidade via divergência estatística de Ng~\cite{ng2011sparse} e a formulação variacional de Kingma e Welling~\cite{kingma2014vae}. Por fim, arquiteturas generativas, baseadas em atenção e em funções \emph{spline} foram embasadas em Goodfellow et al.~\cite{goodfellow2014generative}, Vaswani et al.~\cite{vaswani2017attention} e Liu et al.~\cite{liu2024kan}.

\subsection{3.3. Conceitualização}

Na etapa de conceitualização, adotou-se uma estratégia \emph{middle-out}~\cite{fernandez1997methontology}, partindo dos conceitos centrais mais representativos do aprendizado de máquina para consolidar cinco superclasses de topo mutuamente disjuntas: \onto{Algorithms}, \onto{MLTypes}, \onto{MathematicalFoundations}, \onto{EvaluationMetrics} e \onto{Applications}. A partir desse núcleo, realizou-se simultaneamente a generalização para os termos mais abrangentes e a especialização em subfamílias conceituais, construindo o glossário controlado do domínio e os diagramas conceituais de relacionamento.

Sob essa perspectiva, os algoritmos foram organizados separando-se sua identidade arquitetural do paradigma de aprendizado empregado, permitindo que métodos híbridos ou de fronteira se associem a múltiplos paradigmas de forma consistente. As famílias algorítmicas estruturaram-se conceitualmente em três grandes blocos evolutivos: técnicas de agrupamento (cobrindo abordagens por partição, hierarquia, densidade, grade, dados categóricos, modelos probabilísticos, agrupamento profundo e aprendizado semi-supervisionado)~\cite{nagpal2013review,decarvalho2018gaussian,simoes2023gaussian,yin2026rapid}, algoritmos baseados em árvores (distinguindo modelos individuais de estratégias paralelas e sequenciais em comitê)~\cite{breiman1984cart,quinlan1993c45,freund1997decision,breiman2001random,geurts2006extremely,bentejac2021comparative} e redes neurais artificiais (diferenciando modelos clássicos e profundos, arquiteturas convolucionais, recorrentes, generativas, baseadas em atenção, fundamentadas em \emph{splines} e autocodificadores)~\cite{goodfellow2016deep,rumelhart1986learning,lecun1998gradient,hochreiter1997lstm,vaswani2017attention,liu2024kan}. Paralelamente, organizaram-se as bases matemáticas em áreas clássicas (como cálculo, álgebra linear, otimização, probabilidade, teoria da informação, processamento de sinais e regularização), as métricas conforme a dependência de gabarito e três atributos operacionais referentes à exigência prévia do número de grupos, à robustez a ruído ou valores atípicos e ao sentido de otimização dos indicadores.

\subsection{3.4. Formalização}

Na etapa de formalização, o modelo conceitual foi convertido em uma representação formal computacional baseada na linguagem ontológica padrão da Web Semântica e fundamentada em Lógicas Descritivas decidíveis. As relações taxonômicas foram especificadas por hierarquias estritas de generalização e especialização, complementadas por axiomas explícitos de disjunção entre ramos semanticamente excludentes, preservando a herança múltipla exclusivamente nos casos de modelos híbridos ou conceitos matemáticos transversais.

As conexões interdomínio foram formalizadas mediante propriedades de objeto dotadas de restrições explícitas de domínio, contradomínio, relações inversas, assimetria e irreflexividade. De modo complementar, as características operacionais dos algoritmos e das métricas foram codificadas por meio de propriedades de dados funcionais associadas a valores booleanos e a domínios enumerados controlados. As definições conceituais foram enriquecidas com restrições existenciais, restrições de valor e axiomas de equivalência definidos ao longo dos ciclos evolutivos --- incluindo a caracterização formal de processos baseados em cálculo multivariável e otimização por gradiente --- além da adoção de dualidade entre classes e instâncias para viabilizar consultas analíticas uniformes entre o nível terminológico e o nível assertivo.

\subsection{3.5. Implementação}

A implementação computacional da ontologia foi executada no editor \textbf{Protégé}, permitindo estruturar classes, subclasses, propriedades de objeto, propriedades de dados, restrições semânticas, axiomas de disjunção e indivíduos representativos de algoritmos, fundamentos matemáticos, métricas e aplicações práticas. Todo o vocabulário foi enriquecido com rótulos bilíngues, comentários descritivos e anotações de proveniência bibliográfica visando à rastreabilidade científica e à transparência semântica.

O artefato resultante foi serializado em arquivo no formato padrão RDF/Turtle, consolidando as extensões incrementais dos sucessivos ciclos evolutivos. Durante o processo de codificação, promoveu-se a verificação contínua da integridade estrutural e da ausência de duplicações conceituais, assegurando conformidade com as convenções de modelagem, coesão modular entre os diferentes ramos temáticos e prontidão para inferência automática e recuperação via consultas semânticas.

\subsection{3.6. Validação e Questões de Competência (CQs)}

A validação da ontologia combinou a verificação de consistência lógica com a avaliação funcional orientada a \textbf{Questões de Competência (CQs)}, seguindo práticas consolidadas na engenharia de ontologias~\cite{bezerra2013evaluating,potoniec2020dataset}. No plano lógico-estrutural, a consistência taxonômica, a satisfatibilidade de todas as 213 classes e a coerência dos indivíduos sob herança múltipla foram avaliadas visando ao uso do raciocinador \textbf{HermiT} no Protégé, verificando-se ausência de classes insatisfatíveis, ausência de conflitos nas propriedades de domínio, contradomínio, inversibilidade, assimetria e irreflexividade, bem como plena concordância entre as restrições conceituais e os atributos operacionais preenchidos nas instâncias.

No plano funcional, definiram-se dez Questões de Competência (\textbf{CQ~1} a \textbf{CQ~10}) que atravessam os três ciclos de desenvolvimento, avaliando desde características matemáticas e paradigmas de aprendizado até o uso de atributos booleanos, sentidos ótimos de métricas e aplicabilidade prática. Todas as dez questões foram traduzidas em consultas \textbf{SPARQL} consolidadas no arquivo \texttt{SPARQL\_Consultas.txt}, cujas perguntas e respostas esperadas encontram-se sintetizadas na Tabela~\ref{tab:cqs}.

\begin{table}[htbp]
\centering
\caption{Questões de Competência (CQs) definidas e validadas na OntoML-Mod.}
\label{tab:cqs}
\footnotesize
{\sloppy
\begin{tabularx}{\textwidth}{@{}>{\raggedright\arraybackslash}p{1.05cm}>{\raggedright\arraybackslash}p{5.75cm}>{\raggedright\arraybackslash}X@{}}
\toprule
\textbf{CQ} & \textbf{Pergunta de Competência} & \textbf{Resposta Esperada} \\
\midrule
\textbf{CQ 1} & Quais algoritmos de redes neurais artificiais e de agrupamento profundo utilizam simultaneamente regra da cadeia e descida do gradiente? & Lista de 16 algoritmos: ContractiveAE; Contrastive; ConvolutionalNN; DEC; DenoisingAE; GenerativeAdversarialNet; Kolmogorov-ArnoldNet; LongShort-TermMemoryNet; MultiLayerPerceptron; OverCompleteAE; SparseAE; StackedAE; Transformers; UnderCompleteAE; VaDE; VariationalAE. \\
\addlinespace[3pt]
\textbf{CQ 2} & Quais são os fundamentos matemáticos mais utilizados pelos algoritmos de aprendizado de máquina, em ordem decrescente? & Espera-se uma lista em ordem decrescente dos 44 fundamentos matemáticos por frequência de uso, liderada por PartialDerivatives (31), OptimizationTheories (22), GradientDescent (21), ChainRule (16) e MatrixOperations (15). \\
\addlinespace[3pt]
\textbf{CQ 3} & Qual é a distribuição e a predominância dos paradigmas de aprendizado entre os algoritmos? & UnsupervisedLearning (43); SupervisedLearning (16); SelfSupervisedLearning (10); TransferLearning (3); SemiSupervisedLearning (2). \\
\addlinespace[3pt]
\textbf{CQ 4} & Quais domínios de aplicação são resolvidos por algoritmos da família k-médias? & Análise de Redes; Bioinformática; Mineração de Dados; Reconhecimento de Padrões; Recuperação de Informação; Segmentação de Imagens; Segmentação de Varejo (Clientes). \\
\addlinespace[3pt]
\textbf{CQ 5} & Quais algoritmos empregam conceitos de entropia ou regularização por entropia? & Entropia: C4\_5; COOLCAT; ID3; KFCM\_Ew1; KFCM\_Ew2. Regularização por Entropia: KFCM\_Ew1; KFCM\_Ew2; Kolmogorov-ArnoldNet. \\
\addlinespace[3pt]
\textbf{CQ 6} & Quais algoritmos de agrupamento dispensam a definição prévia do número de grupos e lidam com ruído e valores atípicos? & BIRCH; DBSCAN; HDBSCAN; MeanShift; OPTICS; STING; WaveCluster. \\
\addlinespace[3pt]
\textbf{CQ 7} & Quais métricas de avaliação devem ser maximizadas e quais devem ser minimizadas? & Espera-se uma lista das 23 métricas agrupadas por sentido: 16 para maximizar (ex.: Accuracy, F1Score, Precision, Recall, ROC\_AUC, SilhouetteScore, AdjustedRandIndex) e 7 para minimizar (DaviesBouldinIndex; ErrorOfClassification; Inertia; LogLoss; MAE; MSE; RMSE). \\
\addlinespace[3pt]
\textbf{CQ 8} & Quais modelos baseados em árvores combinam expansão de Taylor de segunda ordem e regularização $L_2$? & CategoricalBoosting; ExtremeGB; LightGBMachine. \\
\addlinespace[3pt]
\textbf{CQ 9} & Quais domínios de aplicação são resolvidos por modelos baseados em árvores e por quais algoritmos? & Espera-se a lista dos 6 domínios (Cibersegurança; Detecção de Anomalias; Detecção de Fraudes; Diagnóstico em Saúde; Mineração de Dados; Reconhecimento de Padrões) pareados aos respectivos algoritmos (CART, RandomForest, ExtremeGB, C4\_5, AdaBoost, LightGBMachine e CategoricalBoosting). \\
\addlinespace[3pt]
\textbf{CQ 10} & Quais domínios de aplicação são resolvidos por arquiteturas de redes neurais artificiais e por quais modelos? & Espera-se a lista dos 8 domínios (Bioinformática; Cibersegurança; Detecção de Anomalias; Diagnóstico em Saúde; Mineração de Dados; Reconhecimento de Padrões; Recuperação de Informação; Segmentação de Imagens) pareados às arquiteturas aplicáveis (como ConvolutionalNN, LongShort-TermMemoryNet, Transformers, GenerativeAdversarialNet, VariationalAE e MultiLayerPerceptron). \\
\bottomrule
\end{tabularx}
}
\end{table}

\subsection{3.7. Modelagem Ontológica}

{\sloppy
A estrutura consolidada da \textbf{OntoML-Mod} compreende \textbf{213 classes} organizadas hierarquicamente sob cinco superclasses de topo mutuamente disjuntas, interligadas por quatro pares de propriedades de objeto inversas e três propriedades de dados funcionais:

\begin{itemize}[leftmargin=*,nosep]
  \item \textbf{1. \onto{Algorithms} (109 classes):} concentra a taxonomia arquitetural dos modelos de aprendizado de máquina em três subárvores principais:
    \begin{itemize}[leftmargin=1.2em,nosep]
      \item \onto{ClusteringAlgorithms} (75 classes), subdividida em oito subfamílias não disjuntas:
        \begin{itemize}[leftmargin=1.2em,nosep]
          \item \onto{PartitionClustering}: inclui \onto{KMeans} (com \onto{HardKMeans}, \onto{FuzzyKMeans}, \onto{SphericalKMeans} e \onto{KernelKMeans}, que abrange \onto{KernelFuzzyKMeans} [\onto{KFCM\_K}, \onto{KFCM\_Ew}, \onto{KFCM\_Ew1}, \onto{KFCM\_Ew2}, \onto{KFCM\_W}, \onto{KFCM\_W1} e \onto{KFCM\_W2}] e \onto{KernelHardKMeans} [\onto{KCM\_F}, \onto{KCM\_F\_GH}, \onto{KCM\_F\_LH}, \onto{KCM\_K}, \onto{KCM\_K\_GH} e \onto{KCM\_K\_LH}]), \onto{KMedoids} (\onto{CLARA}, \onto{CLARANS}, \onto{PAM}), \onto{KModes} e \onto{KPrototype};
          \item \onto{HierarchicalClustering}: \onto{Agglomerative}, \onto{Divisive}, \onto{SingleLinkage}, \onto{CompleteLinkage}, \onto{AverageLinkage}, \onto{BIRCH}, \onto{CURE}, \onto{CHAMELEON}, \onto{CACTUS}, \onto{ROCK} e \onto{LIMBO};
          \item \onto{DensityClustering}: \onto{DBSCAN}, \onto{HDBSCAN}, \onto{OPTICS}, \onto{MeanShift}, \onto{DENCLUE}, \onto{DBCLASD}, \onto{CLIQUE} e \onto{WaveCluster};
          \item \onto{GridClustering}: \onto{STING}, \onto{CLIQUE}, \onto{WaveCluster}, \onto{GridDBSCAN} e \onto{STDBSCAN};
          \item \onto{CategoricalClustering}: \onto{CACTUS}, \onto{COOLCAT}, \onto{KModes}, \onto{KPrototype}, \onto{LIMBO}, \onto{ROCK} e \onto{STIRR};
          \item \onto{ModelClustering}: \onto{AutoClass}, \onto{CLASSIT}, \onto{COBWEB}, \onto{ExpectationMaximization}, \onto{GaussianMixtureModels}, \onto{HiddenMarkovModels}, \onto{LatentDirichletAllocation} e \onto{Self-OrganizingMap};
          \item \onto{DeepLearningClustering}: \onto{DEC}, \onto{IDEC}, \onto{VaDE}, \onto{DCEC}, \onto{DCN} e \onto{Contrastive};
          \item \onto{SemiSupervisedClustering}: \onto{ConstrainedKMeans} e \onto{SemiSupervisedFuzzyCMeans}.
        \end{itemize}
      \item \onto{TreeBased} (16 classes), organizada em três ramos disjuntos:
        \begin{itemize}[leftmargin=1.2em,nosep]
          \item \onto{DecisionTree}: \onto{ID3}, \onto{C4\_5} e \onto{CART};
          \item \onto{ParallelTreeEnsemble}: \onto{ExtraTrees} e \onto{TreeBagging} (que especializa em \onto{RandomForest});
          \item \onto{TreeBoosting}: \onto{AdaBoost} e \onto{TreeGradientBoosting} (com \onto{GradientBoostingMachine}, \onto{ExtremeGB}, \onto{LightGBMachine} e \onto{CategoricalBoosting}).
        \end{itemize}
      \item \onto{ArtificialNeuralNetwork} (18 classes), integrando redes clássicas e profundas:
        \begin{itemize}[leftmargin=1.2em,nosep]
          \item Redes clássicas e mapas auto-organizáveis: \onto{MultiLayerPerceptron}, \onto{RadialBasisFunctionNet} e \onto{Self-OrganizingMap} (esta última compartilhada por herança múltipla com \onto{ModelClustering}~\cite{nagpal2013review});
          \item \onto{DeepNeuralNetwork} (14 classes): \onto{ConvolutionalNN}, \onto{LongShort-TermMemoryNet}, \onto{GenerativeAdversarialNet}, \onto{Transformers}, \onto{Kolmogorov-ArnoldNet} e a subfamília \onto{AutoEncoders} (\onto{ContractiveAE}, \onto{DenoisingAE}, \onto{OverCompleteAE}, \onto{SparseAE}, \onto{StackedAE}, \onto{UnderCompleteAE} e \onto{VariationalAE}).
        \end{itemize}
    \end{itemize}

  \item \textbf{2. \onto{MathematicalFoundations} (53 classes):} estruturada em nove subáreas disjuntas:
    \begin{itemize}[leftmargin=1.2em,nosep]
      \item \onto{Calculus}: \onto{ActivationFunctions}, \onto{MultivariableCalculus} (\onto{ChainRule}, \onto{JacobianMatrix}, \onto{PartialDerivatives}), \onto{SplineFunctions} e \onto{TaylorExpansion};
      \item \onto{DistanceAndSimilarityMetrics}: \onto{CosineSimilarity}, \onto{EuclideanDistance}, \onto{JaccardCoefficient}, \onto{ManhattanDistance} e \onto{KernelFunctions} (\onto{GaussianKernel});
      \item \onto{GraphTheory};
      \item \onto{InformationTheory}: \onto{CrossEntropy}, \onto{Entropy}, \onto{GainRatio}, \onto{InformationGain} e \onto{KullbackLeiblerDivergence};
      \item \onto{LinearAlgebra}: \onto{AttentionMechanism}, \onto{EigenvaluesAndEigenvectors} e \onto{MatrixOperations};
      \item \onto{OptimizationTheories}: \onto{GradientDescent}, \onto{LagrangeMultipliers} e \onto{MinimaxOptimization};
      \item \onto{ProbabilityAndStatistics}: \onto{BayesianInference}, \onto{BootstrapResampling}, \onto{ImpurityMeasures} (\onto{GiniImpurity}, \onto{VarianceReduction}) e \onto{ProbabilityDistributions} (\onto{GaussianDistribution});
      \item \onto{RegularizationTechniques}: \onto{BatchNormalization}, \onto{DataAugmentation}, \onto{Dropout}, \onto{EarlyStopping}, \onto{ElasticNet}, \onto{EntropyRegularization}, \onto{TreePruning} e \onto{WeightDecay} (\onto{L1\_Reg}, \onto{L2\_Reg});
      \item \onto{SignalProcessing}: \onto{ConvolutionOperation} e \onto{WaveletTransform}.
    \end{itemize}

  \item \textbf{3. \onto{EvaluationMetrics} (31 classes):} dividida segundo a necessidade de rótulos de referência:
    \begin{itemize}[leftmargin=1.2em,nosep]
      \item \onto{GroundTruthDependentMetrics}: \onto{ClassificationMetrics} (\onto{Accuracy}, \onto{ErrorOfClassification}, \onto{F1Score}, \onto{LogLoss}, \onto{Precision}, \onto{ROC\_AUC} e \onto{Recall}), \onto{ClusterMatchingMetrics} (\onto{AdjustedRandIndex}, \onto{FowlkesMallowsIndex}, \onto{FuzzyRandIndexFrigui}, \onto{FuzzyRandIndexHullermeier}, \onto{JaccardIndex} e \onto{NormalizedMutualInformation}) e \onto{RegressionMetrics} (\onto{MAE}, \onto{MSE} e \onto{RMSE});
      \item \onto{GroundTruthIndependentMetrics}: \onto{ClusterValidityIndices} (\onto{BIC}, \onto{Calinski-Harabasz}, \onto{DaviesBouldinIndex}, \onto{DunnsIndex}, \onto{GapStatistic}, \onto{Inertia}, \onto{ModifiedPartitionCoefficient} e \onto{SilhouetteScore}).
    \end{itemize}

  \item \textbf{4. \onto{Applications} (13 classes):} reúne os domínios de aplicação \onto{AnomalyDetection}, \onto{Bioinformatics}, \onto{ChurnPrediction} (reservada para expansão futura), \onto{CustomerSegmentation}, \onto{Cybersecurity}, \onto{DataMining}, \onto{FraudDetection}, \onto{HealthcareDiagnosis}, \onto{ImageSegmentation}, \onto{InformationRetrieval}, \onto{NetworkAnalysis} e \onto{PatternRecognition}, povoados por onze instâncias nominais interligadas aos algoritmos.

  \item \textbf{5. \onto{MLTypes} (7 classes):} particiona os paradigmas em \onto{SupervisedLearning}, \onto{UnsupervisedLearning}, \onto{SemiSupervisedLearning}, \onto{SelfSupervisedLearning}, \onto{ReinforcementLearning} e \onto{TransferLearning}.
\end{itemize}

A integração semântica entre esses cinco eixos ocorre por meio de quatro pares de propriedades de objeto inversas --- \onto{usesMathematicalFoundation} e \onto{IsMathematicalUsedBy}; \onto{utilizesMLType} e \onto{isMLTypeOf}; \onto{hasEvaluationMetric} e \onto{isMetricOf}; e \onto{isUsedInApplication} e \onto{isSolvedBy} --- complementados pelas três propriedades de dados funcionais \onto{requiresClusterCountK}, \onto{handlesNoiseOutliers} e \onto{hasOptimalDirection}. Dessa forma, cada algoritmo herda o perfil de sua subárvore e acrescenta assinaturas conceituais próprias, como exemplificado por \onto{ExtremeGB} (que associa árvores com gradiente descendente à expansão de Taylor, regularização $L_1$/$L_2$ e poda) e pelas definições equivalentes introduzidas no Ciclo~3 para caracterizar famílias neurais --- distinguindo, por exemplo, \onto{Transformers} pelo uso de mecanismos de atenção, \onto{GenerativeAdversarialNet} pela otimização minimax, \onto{Kolmogorov-ArnoldNet} por funções \emph{spline} e \onto{VariationalAE} pela combinação de inferência bayesiana e divergência de Kullback-Leibler sobre autocodificadores.
}

% ==============================================================================
% Seção 7 numerada estritamente conforme template_de_prompt_para_artigo.md
% (Seções 4, 5 e 6 reservadas para complementação posterior do artigo)
% ==============================================================================
\renewcommand{\refname}{7. REFERÊNCIAS BIBLIOGRÁFICAS}
{\sloppy
\begin{thebibliography}{99}\small

\bibitem{baldi2012autoencoders}
Baldi, P. Autoencoders, Unsupervised Learning, and Deep Architectures. In: \emph{Proceedings of ICML Workshop on Unsupervised and Transfer Learning}, PMLR, v. 27, p. 37--49, 2012.

\bibitem{bengio2007greedy}
Bengio, Y.; Lamblin, P.; Popovici, D.; Larochelle, H. Greedy Layer-Wise Training of Deep Networks. In: \emph{Advances in Neural Information Processing Systems 19 (NIPS 2006)}, p. 153--160, 2007.

\bibitem{bentejac2021comparative}
Bentéjac, C.; Csörgő, A.; Martínez-Muñoz, G. A comparative analysis of gradient boosting algorithms. \emph{Artificial Intelligence Review}, v. 54, n. 3, p. 1937--1967, 2021. DOI: \texttt{10.1007/s10462-020-09896-5}.

\bibitem{bezerra2013evaluating}
Bezerra, C.; Freitas, F.; Santana, F. Evaluating Ontologies with Competency Questions. In: \emph{2013 IEEE/WIC/ACM International Joint Conferences on Web Intelligence (WI) and Intelligent Agent Technologies (IAT)}, p. 284--285, 2013. DOI: \texttt{10.1109/WI-IAT.2013.199}.

\bibitem{braga2020mlonto}
Braga, J.; Dias, J. L. R.; Regateiro, F. A Machine Learning Ontology (MLOnto). Preprint, Open Science Framework (OSF), Nov. 2020. DOI: \texttt{10.17605/OSF.IO/CHU5Q}.

\bibitem{breiman1984cart}
Breiman, L.; Friedman, J. H.; Olshen, R. A.; Stone, C. J. \emph{Classification and Regression Trees}. Wadsworth / CRC Press, 1984. 368 p. ISBN \texttt{978-0412048418}.

\bibitem{breiman2001random}
Breiman, L. Random Forests. \emph{Machine Learning}, v. 45, n. 1, p. 5--32, 2001. DOI: \texttt{10.1023/A:1010933404324}.

\bibitem{broomhead1988multivariable}
Broomhead, D. S.; Lowe, D. Multivariable Functional Interpolation and Adaptive Networks. \emph{Complex Systems}, v. 2, n. 3, p. 321--355, 1988.

\bibitem{decarvalho2018gaussian}
de Carvalho, F. A. T.; Simões, E. C.; Santana, L. V. C.; Ferreira, M. R. P. Gaussian kernel c-means hard clustering algorithms with automated computation of the width hyper-parameters. \emph{Pattern Recognition}, v. 79, p. 370--386, 2018. DOI: \texttt{10.1016/j.patcog.2018.02.018}.

\bibitem{dhillon2001concept}
Dhillon, I. S.; Modha, D. S. Concept Decompositions for Large Sparse Text Data Using Clustering. \emph{Machine Learning}, v. 42, n. 1--2, p. 143--175, 2001.

\bibitem{elmassari2022diabetes}
El~Massari, H.; Sabouri, Z.; Mhammedi, S.; Gherabi, N. Diabetes Prediction Using Machine Learning Algorithms and Ontology. \emph{Journal of ICT Standardization}, v. 10, n. 2, p. 319--338, 2022. DOI: \texttt{10.13052/jicts2245-800X.10212}.

\bibitem{fernandez1997methontology}
Fernández, M.; Gómez-Pérez, A.; Juristo, N. METHONTOLOGY: From Ontological Art Towards Ontological Engineering. In: \emph{Proceedings of the AAAI97 Spring Symposium on Ontological Engineering}, AAAI Technical Report SS-97-06, p. 33--40, 1997.

\bibitem{freund1997decision}
Freund, Y.; Schapire, R. E. A Decision-Theoretic Generalization of On-Line Learning and an Application to Boosting. \emph{Journal of Computer and System Sciences}, v. 55, n. 1, p. 119--139, 1997. DOI: \texttt{10.1006/jcss.1997.1504}.

\bibitem{geurts2006extremely}
Geurts, P.; Ernst, D.; Wehenkel, L. Extremely Randomized Trees. \emph{Machine Learning}, v. 63, n. 1, p. 3--42, 2006. DOI: \texttt{10.1007/s10994-006-6226-1}.

\bibitem{gibson1998clustering}
Gibson, D.; Kleinberg, J.; Raghavan, P. Clustering Categorical Data: An Approach Based on Dynamical Systems. In: \emph{Proceedings of the 24th International Conference on Very Large Data Bases (VLDB)}, 1998.

\bibitem{goodfellow2014generative}
Goodfellow, I. J.; Pouget-Abadie, J.; Mirza, M.; Xu, B.; Warde-Farley, D.; Ozair, S.; Courville, A.; Bengio, Y. Generative Adversarial Nets. In: \emph{Advances in Neural Information Processing Systems 27 (NIPS 2014)}, v. 2, p. 2672--2680, 2014.

\bibitem{goodfellow2016deep}
Goodfellow, I.; Bengio, Y.; Courville, A. \emph{Deep Learning}. MIT Press, 2016. 800 p. ISBN \texttt{978-0262035618}.

\bibitem{hochreiter1997lstm}
Hochreiter, S.; Schmidhuber, J. Long Short-Term Memory. \emph{Neural Computation}, v. 9, n. 8, p. 1735--1780, 1997. DOI: \texttt{10.1162/neco.1997.9.8.1735}.

\bibitem{humm2026industry}
Humm, B. G. An Industry-Ready Machine Learning Ontology. \emph{Applied Sciences}, v. 16, n. 2, 843, 2026. DOI: \texttt{10.3390/app16020843}.

\bibitem{kingma2014vae}
Kingma, D. P.; Welling, M. Auto-Encoding Variational Bayes. In: \emph{2nd International Conference on Learning Representations (ICLR 2014)}, \texttt{arXiv:1312.6114}, 2014.

\bibitem{lecun1998gradient}
LeCun, Y.; Bottou, L.; Bengio, Y.; Haffner, P. Gradient-Based Learning Applied to Document Recognition. \emph{Proceedings of the IEEE}, v. 86, n. 11, p. 2278--2324, 1998. DOI: \texttt{10.1109/5.726791}.

\bibitem{liu2023self}
Liu, X.; Zhang, F.; Hou, Z.; Mian, L.; Wang, Z.; Zhang, J.; Tang, J. Self-Supervised Learning: Generative or Contrastive. \emph{IEEE Transactions on Knowledge and Data Engineering}, v. 35, n. 1, p. 857--876, 2023 (versão local consultada: \texttt{arXiv:2006.08218v5}, Mar. 2021). DOI: \texttt{10.1109/TKDE.2021.3090866}.

\bibitem{liu2024kan}
Liu, Z.; Wang, Y.; Vaidya, S.; Ruehle, F.; Halverson, J.; Soljačić, M.; Hou, T. Y.; Tegmark, M. KAN: Kolmogorov-Arnold Networks. Preprint \texttt{arXiv:2404.19756 [cs.LG]}, 2024.

\bibitem{nagpal2013review}
Nagpal, A.; Jatain, A.; Gaur, D. Review based on data clustering algorithms. In: \emph{2013 IEEE Conference on Information \& Communication Technologies (ICT)}, p. 298--303, 2013. DOI: \texttt{10.1109/CICT.2013.6558109}.

\bibitem{ng2011sparse}
Ng, A. Sparse Autoencoder. \emph{CS294A Lecture Notes}, Stanford University, 2011.

\bibitem{potoniec2020dataset}
Potoniec, J.; Wiśniewski, D.; Ławrynowicz, A.; Keet, C. M. Dataset of ontology competency questions to SPARQL-OWL queries translations. \emph{Data in Brief}, v. 29, 105098, 2020. DOI: \texttt{10.1016/j.dib.2019.105098}.

\bibitem{publio2018mlschema}
Publio, G.; Ławrynowicz, A.; Panov, P.; Soldatova, L.; Soru, T.; Vanschoren, J.; Zafar, H. ML-Schema: Exposing the Semantics of Machine Learning with Schemas and Ontologies. Preprint \texttt{arXiv:1807.05351v1 [cs.LG]}, Jul. 2018.

\bibitem{quinlan1993c45}
Quinlan, J. R. \emph{C4.5: Programs for Machine Learning}. Morgan Kaufmann, 1993. 302 p. ISBN \texttt{978-1558602380}.

\bibitem{rifai2011contractive}
Rifai, S.; Vincent, P.; Muller, X.; Glorot, X.; Bengio, Y. Contractive Auto-Encoders: Explicit Invariance During Feature Extraction. In: \emph{Proceedings of the 28th International Conference on Machine Learning (ICML 2011)}, p. 833--840, 2011.

\bibitem{rumelhart1986learning}
Rumelhart, D. E.; Hinton, G. E.; Williams, R. J. Learning representations by back-propagating errors. \emph{Nature}, v. 323, p. 533--536, 1986. DOI: \texttt{10.1038/323533a0}.

\bibitem{sarker2021ml}
Sarker, I. H. Machine Learning: Algorithms, Real-World Applications and Research Directions. \emph{SN Computer Science}, v. 2, n. 3, 160, 2021. DOI: \texttt{10.1007/s42979-021-00592-x}.

\bibitem{simoes2023gaussian}
Simões, E. C.; de Carvalho, F. A. T. Gaussian kernel fuzzy c-means with width parameter computation and regularization. \emph{Pattern Recognition}, v. 143, 109749, 2023. DOI: \texttt{10.1016/j.patcog.2023.109749}.

\bibitem{vaswani2017attention}
Vaswani, A.; Shazeer, N.; Parmar, N.; Uszkoreit, J.; Jones, L.; Gomez, A. N.; Kaiser, Ł.; Polosukhin, I. Attention Is All You Need. In: \emph{Advances in Neural Information Processing Systems 30 (NIPS 2017)}, v. 30, p. 5998--6008, 2017.

\bibitem{vincent2008extracting}
Vincent, P.; Larochelle, H.; Bengio, Y.; Manzagol, P.-A. Extracting and composing robust features with denoising autoencoders. In: \emph{Proceedings of the 25th International Conference on Machine Learning (ICML 2008)}, p. 1096--1103, 2008. DOI: \texttt{10.1145/1390156.1390294}.

\bibitem{vincent2010stacked}
Vincent, P.; Larochelle, H.; Lajoie, I.; Bengio, Y.; Manzagol, P.-A. Stacked Denoising Autoencoders: Learning Useful Representations in a Deep Network with a Local Denoising Criterion. \emph{Journal of Machine Learning Research}, v. 11, p. 3371--3408, 2010.

\bibitem{yin2026rapid}
Yin, H.; Aryani, A.; Petrie, S.; Nambissan, A.; Astudillo, A.; Cao, S. A rapid review of clustering algorithms. \emph{Array}, v. 30, 100904, 2026. DOI: \texttt{10.1016/j.array.2026.100904}.

\bibitem{zhuang2021transfer}
Zhuang, F.; Qi, Z.; Duan, K.; Xi, D.; Zhu, Y.; Zhu, H.; Xiong, H.; He, Q. A Comprehensive Survey on Transfer Learning. \emph{Proceedings of the IEEE}, v. 109, n. 1, p. 43--76, 2021 (versão local consultada: \texttt{arXiv:1911.02685v3}, Jun. 2020). DOI: \texttt{10.1109/JPROC.2020.3004555}.

\end{thebibliography}
}

\end{document}
